%% toc: P15 160 160.9
% Proofs for Feller's Poisson, Laplace and analytic moment continuations (Feller II, VII.5--VII.6 and XV.4).
\section{Feller's Poisson, Laplace and analytic moment continuations}
\label{sec:annex-p15-w02b-p15-012-feller-continuations}

\begin{annexitem}{annex:p15-w02b-p15-012-poisson-identity}{thm:p15-w02b-p15-012-poisson-identity}{The Poisson identity by absolute rearrangement}
\begin{proof}
Let \(|a_j|\leq A\) for all \(j\). Expanding the iterated difference gives
\[
 \Delta^r a_i=\sum_{k=0}^r\binom rk(-1)^{r-k}a_{i+k},
\]
which is the case \(r=0\) trivially and follows for \(r+1\) from the case
\(r\) by \(\Delta^{r+1}a_i=\Delta^ra_{i+1}-\Delta^ra_i\) and Pascal's rule.
Consider the double series
\[
 \sum_{r=0}^\infty\sum_{k=0}^r\frac{x^r}{r!}\binom rk(-1)^{r-k}a_{i+k}.
\]
The sum of the absolute values of its terms is at most
\(\sum_r|x|^r2^rA/r!=Ae^{2|x|}<\infty\). Hence the double series converges
absolutely, its value equals the left side of (5.1), and its terms may be
summed in any order. Put \(r=k+m\) with \(k,m\geq0\). Since
\(\binom{k+m}k/(k+m)!=1/(k!\,m!)\), the sum becomes
\[
 \sum_{k=0}^\infty\frac{x^k}{k!}a_{i+k}\sum_{m=0}^\infty\frac{(-x)^m}{m!}
     =e^{-x}\sum_{k=0}^\infty\frac{x^k}{k!}a_{i+k},
\]
and the last series converges absolutely because \(|a_{i+k}|\leq A\). This is
(5.1).

For (5.2), put \(a_j=u(jh)\), \(i=0\) and \(x=\theta/h\). By the definition of
the normalized difference, \(\Delta_h^ru(0)=h^{-r}\Delta^ra_0\), so the term
\(x^r\Delta^ra_0/r!\) equals \(\theta^r\Delta_h^ru(0)/r!\). The right side of
(5.1) becomes \(e^{-\theta/h}\sum_ju(jh)(\theta/h)^j/j!\), which is
\(\mathbb E\,u(hN_{\theta/h})\). The variable \(hN_{\theta/h}\) has mean
\(h\cdot\theta/h=\theta\) and variance \(h^2\cdot\theta/h=h\theta\).
\end{proof}
\end{annexitem}

\begin{annexitem}{annex:p15-w02b-p15-012-generalized-taylor}{thm:p15-w02b-p15-012-generalized-taylor}{The generalized Taylor formula and the two examples}
\begin{proof}
Fix \(t\geq0\) and put \(v(x)=u(t+x)\) for \(x\geq0\). Then \(v\) is bounded
and continuous with \(\sup|v|\leq\|u\|\), and \(\Delta_h^rv(0)=\Delta_h^ru(t)\)
for every \(r\), since both are built from the same values
\(u(t+jh)\). By \cref{thm:p15-w02b-p15-012-poisson-identity},
\[
 \sum_{r=0}^\infty\frac{\theta^r}{r!}\Delta_h^ru(t)=\mathbb E\,v(hN_{\theta/h}).
\]
At \(\theta=0\) the variable \(hN_0\) is zero, so the value is \(v(0)=u(t)\).

Let \(b>0\) and \(\epsilon>0\). The function \(v\) is uniformly continuous on
\([0,b+1]\), so there is \(\delta\in(0,1]\) with \(|v(x)-v(y)|<\epsilon\)
whenever \(x,y\in[0,b+1]\) and \(|x-y|<\delta\). Let
\(0\leq\theta\leq b\) and write \(Y=hN_{\theta/h}\), a nonnegative variable.
On the event \(|Y-\theta|<\delta\) we have \(Y\in[0,b+1]\), so
\(|v(Y)-v(\theta)|<\epsilon\). On the complementary event
\(|v(Y)-v(\theta)|\leq2\|u\|\), and by Chebyshev's inequality that event has
probability at most \(\operatorname{Var}(Y)/\delta^2=h\theta/\delta^2\leq
hb/\delta^2\). Therefore
\[
 \bigl|\mathbb E\,v(Y)-v(\theta)\bigr|\leq\epsilon+\frac{2\|u\|hb}{\delta^2},
\]
and the right side is below \(2\epsilon\) for all \(0\leq\theta\leq b\) once
\(h<\epsilon\delta^2/(2\|u\|b+1)\). Since \(v(\theta)=u(t+\theta)\) and
\(\epsilon\) is arbitrary, (5.3) holds uniformly for \(\theta\in[0,b]\).

If \(u\) is uniformly continuous on \([0,\infty)\), choose \(\delta\) so that
\(|u(x)-u(y)|<\epsilon\) whenever \(|x-y|<\delta\). This \(\delta\) serves
for every \(v\), that is, for every \(t\geq0\), and the same bound gives
convergence uniform in \(t\geq0\) and \(0\leq\theta\leq b\).

For the negative case let \(u\) be bounded and continuous on the line,
\(\theta<0\) and \(h<0\). Put \(a_j=u(t+jh)\) and \(x=\theta/h>0\). Then
\(\Delta^ra_0=h^r\Delta_h^ru(t)\) with the quotient taken at step \(h\), and
(5.1) gives
\[
 \sum_{r=0}^\infty\frac{\theta^r}{r!}\Delta_h^ru(t)
   =\mathbb E\,u(t+hN_{\theta/h}).
\]
The variable \(Y=hN_{\theta/h}\) has mean \(\theta\) and variance
\(h^2\theta/h=h\theta>0\), which tends to zero as \(h\to0-\). The preceding
estimate, with uniform continuity of \(u\) on \([t+\theta-1,t+\theta+1]\),
gives \(\mathbb E\,u(t+Y)\to u(t+\theta)\).

For the first example, \(u(x)=e^{-x}\) gives
\(\Delta_hu(x)=e^{-x}(e^{-h}-1)/h=c_hu(x)\), hence \(\Delta_h^ru=c_h^ru\) by
induction, and the series equals
\(e^{-t}\sum_r(\theta c_h)^r/r!=e^{-t}e^{\theta c_h}\). For \(h>0\),
\(e^{-h}>1-h\), so \(c_h>-1\) and \(e^{-t}e^{\theta c_h}>e^{-t-\theta}\) when
\(\theta>0\). As \(h\to0+\), \(c_h\to-1\), the derivative of \(e^{-x}\) at
\(0\), so the value tends to \(e^{-t-\theta}\).
For the second example, \(u(x)=\min(x,1)\) is bounded and continuous on
\([0,\infty)\). The identity above at \(t=0\) gives the expectation
\(\mathbb E\min(hN_{\theta/h},1)\), and the theorem gives the limit
\(\min(\theta,1)\).
\end{proof}
\end{annexitem}

\begin{annexitem}{annex:p15-w02b-p15-012-exponential-formula}{thm:p15-w02b-p15-012-exponential-formula}{The exponential formula for translations and its operator-norm boundary}
\begin{proof}
For \(u\in C_b\), \(\|T(h)u\|\leq\|u\|\), so \(\|\Delta_hu\|\leq2\|u\|/h\)
and by induction \(\|\Delta_h^ru\|\leq(2/h)^r\|u\|\). The series
\(\sum_r\theta^r\Delta_h^ru/r!\) therefore converges uniformly on
\([0,\infty)\), with sum of norms at most \(e^{2\theta/h}\|u\|\), and each
partial sum is continuous, so the sum lies in \(C_b\). Similarly the
weights of (5.7) sum to \(1\), so \(A_h(\theta)u\) is a uniform limit of
continuous functions, lies in \(C_b\), and has norm at most \(\|u\|\).

Fix \(t\geq0\) and put \(a_j=u(t+jh)\), \(x=\theta/h\). Then
\(\Delta^ra_0=h^r\Delta_h^ru(t)\), and (5.1) reads
\[
 \sum_{r=0}^\infty\frac{\theta^r}{r!}\Delta_h^ru(t)
   =e^{-\theta/h}\sum_{k=0}^\infty\frac1{k!}\Bigl(\frac\theta h\Bigr)^k u(t+kh)
   =A_h(\theta)u(t).
\]
As this holds for every \(u\) and \(t\), the two operators coincide. The
convergence (5.6) at each \(t\) is
\cref{thm:p15-w02b-p15-012-generalized-taylor}, and when \(u\) is uniformly
continuous that theorem gives convergence uniform in \(t\), which is
convergence in the norm of \(C_b\).

For the operator norm, \(\|A_h(\theta)u-T(\theta)u\|\leq2\|u\|\) gives
\(\|A_h(\theta)-T(\theta)\|\leq2\). Let \(\theta>0\) and let \(h>0\) with
\(\theta/h\) not an integer. The grid points \(kh\), \(k\geq0\), then avoid
\(\theta\), and the nearest of them lies at the positive distance
\(d=\min\{\theta-\lfloor\theta/h\rfloor h,\ (\lfloor\theta/h\rfloor+1)h-\theta\}\).
Define \(u(x)=-1+2\max\{0,1-|x-\theta|/d\}\). It is continuous with
\(\|u\|=1\), \(u(\theta)=1\), and \(u(kh)=-1\) for every \(k\geq0\). At
\(t=0\),
\[
 A_h(\theta)u(0)-T(\theta)u(0)=\sum_{k\geq0}e^{-\theta/h}\frac{(\theta/h)^k}{k!}(-1)-1=-2 .
\]
Thus \(\|A_h(\theta)-T(\theta)\|\geq2\), and equality holds. Since such \(h\)
exist in every interval \((0,h_0)\), \(A_h(\theta)\) does not converge to
\(T(\theta)\) in operator norm as \(h\to0+\).
\end{proof}
\end{annexitem}

\begin{annexitem}{annex:p15-w02b-p15-012-laplace-derivatives}{thm:p15-w02b-p15-012-laplace-derivatives}{Differentiation of the Laplace transform and of the Laplace integral}
\begin{proof}
For \(m\geq1\) and \(a>0\), the function \(y\mapsto y^me^{-ay}\) on
\([0,\infty)\) has its maximum at \(y=m/a\), so
\(y^me^{-ay}\leq(m/(ae))^m\).

Fix \(\lambda>0\) and \(\lambda_1\in(0,\lambda)\). Assume (6.3) for some
\(k\geq0\); the case \(k=0\) is the definition. For real \(\eta\neq0\) with
\(|\eta|<\lambda-\lambda_1\), the mean value theorem gives a point \(\xi\)
between \(\lambda\) and \(\lambda+\eta\), depending on \(\theta\), with
\[
 \frac{e^{-(\lambda+\eta)\theta}-e^{-\lambda\theta}}\eta\,\theta^k
   =-\theta^{k+1}e^{-\xi\theta}.
\]
Since \(\xi\geq\lambda_1\), the absolute value is at most
\(\theta^{k+1}e^{-\lambda_1\theta}\leq((k+1)/(\lambda_1e))^{k+1}\), a
constant, which is \(F\)-integrable. As \(\eta\to0\) along any sequence the
quotient tends to \(-\theta^{k+1}e^{-\lambda\theta}\). Dominated convergence
shows that \((-1)^k\varphi^{(k)}\) is differentiable at \(\lambda\) with
derivative \(-\int e^{-\lambda\theta}\theta^{k+1}F\{d\theta\}\), which is
(6.3) for \(k+1\).

For \(w\) the same quotient argument applies with \(F\{d\theta\}\) replaced
by \(u(x)\,dx\) and the dominating function
\(\|u\|x^{k+1}e^{-\lambda_1x}\). This function is Lebesgue integrable on
\([0,\infty)\), since \(x^{k+1}e^{-\lambda_1x/2}\leq(2(k+1)/(\lambda_1e))^{k+1}\)
and \(e^{-\lambda_1x/2}\) is integrable. The base case is
\(w^{(0)}=w\), and the induction gives
\(w^{(k)}(\lambda)=\int_0^\infty(-x)^ke^{-\lambda x}u(x)\,dx\).
\end{proof}
\end{annexitem}

\begin{annexitem}{annex:p15-w02b-p15-012-laplace-inversion}{thm:p15-w02b-p15-012-laplace-inversion}{The Poisson inversion of Laplace transforms}
\begin{proof}
First (6.1). The equality of the two left expressions is the definition of
the Poisson law. If \(x<0\) and \(\theta\geq0\), then \(\theta>x\) and
\(\mathbb P\{N_{\lambda\theta}\leq\lambda x\}=0\). If \(\theta=0<x\), then
\(N_0=0\leq\lambda x\) and the probability is \(1\). Otherwise
\(\theta>0\), and Chebyshev's inequality with mean and variance
\(\lambda\theta\) gives, for \(\theta>x\),
\[
 \mathbb P\{N_{\lambda\theta}\leq\lambda x\}
   \leq\mathbb P\{|N_{\lambda\theta}-\lambda\theta|\geq\lambda(\theta-x)\}
   \leq\frac{\theta}{\lambda(\theta-x)^2}\longrightarrow0,
\]
and for \(\theta<x\) the same bound with \(x-\theta\) controls
\(\mathbb P\{N_{\lambda\theta}>\lambda x\}\), so the probability tends to
\(1\).

Now (6.4). Let \(\lambda>0\). By (6.3) and linearity of the integral over the
finite set \(k\leq\lambda x\),
\[
 \sum_{k\leq\lambda x}\frac{(-1)^k}{k!}\lambda^k\varphi^{(k)}(\lambda)
   =\int_0^\infty e^{-\lambda\theta}\sum_{k\leq\lambda x}\frac{(\lambda\theta)^k}{k!}\,F\{d\theta\}
   =\int_0^\infty\mathbb P\{N_{\lambda\theta}\leq\lambda x\}\,F\{d\theta\}.
\]
If \(x<0\) both sides are zero, as is \(F(x)\). Let \(x\geq0\) be a continuity
point of \(F\), so \(F\{\{x\}\}=0\). The integrand lies in \([0,1]\) and, by
(6.1), tends to \(1\) for \(\theta<x\) and to \(0\) for \(\theta>x\), that
is, for \(F\)-almost every \(\theta\). Dominated convergence gives the limit
\(F\{[0,x)\}=F\{[0,x]\}=F(x)\).

For uniqueness, let \(F\) and \(G\) on \([0,\infty)\) have Laplace transforms
that agree for \(\lambda>\lambda_1\). Their derivatives of every order then
agree there, so the left sides of (6.4) agree for large \(\lambda\), and
\(F(x)=G(x)\) at every common continuity point. A monotone function has at
most countably many discontinuities, so common continuity points are dense.
Both functions are right continuous, so \(F(x)=G(x)\) for every \(x\).
\end{proof}
\end{annexitem}

\begin{annexitem}{annex:p15-w02b-p15-012-laplace-integral-inversion}{thm:p15-w02b-p15-012-laplace-integral-inversion}{Inversion of Laplace integrals by the gamma kernel}
\begin{proof}
By \cref{thm:p15-w02b-p15-012-laplace-derivatives} at \(\lambda=n/\theta>0\),
\[
 \frac{(-1)^{n-1}}{(n-1)!}\Bigl(\frac n\theta\Bigr)^nw^{(n-1)}(n/\theta)
   =\frac{1}{(n-1)!}\Bigl(\frac n\theta\Bigr)^n\int_0^\infty x^{n-1}e^{-nx/\theta}u(x)\,dx,
\]
because \((-1)^{n-1}(-x)^{n-1}=x^{n-1}\). Writing
\((n/\theta)^nx^{n-1}=(nx/\theta)^{n-1}(n/\theta)\) gives the middle expression
of (6.6). It is \(\int u\,d\nu_{n,\theta}\) for the gamma law
\(\nu_{n,\theta}\) of shape \(n\) and rate \(n/\theta\), normalized by
\(\Gamma(n)=(n-1)!\). By \cref{thm:p15-feller-kernel-examples}, with
\(\nu_{n,0}=\delta_0\), this expectation converges to \(u(\theta)\) uniformly
on every compact interval \([0,b]\). Restricting to \(0<\theta\leq b\), where
the left side of (6.6) is defined, gives the assertion.
\end{proof}
\end{annexitem}

\begin{annexitem}{annex:p15-w02b-p15-012-laplace-moment-inequalities}{thm:p15-w02b-p15-012-laplace-moment-inequalities}{The alternating inequalities and the moment series}
\begin{proof}
For \(m\geq0\) put \(R_m(t)=e^{-t}-\sum_{k=0}^m(-t)^k/k!\), and put
\(R_{-1}(t)=e^{-t}\). Then \(R_m(0)=0\) for \(m\geq0\) and
\[
 R_m'(t)=-e^{-t}-\sum_{k=1}^m\frac{(-1)^kt^{k-1}}{(k-1)!}
       =-e^{-t}+\sum_{j=0}^{m-1}\frac{(-t)^j}{j!}=-R_{m-1}(t).
\]
We show \((-1)^{m+1}R_m(t)>0\) for \(t>0\) and \(m\geq-1\). For \(m=-1\) it
is \(e^{-t}>0\). If it holds for \(m-1\), then for \(t>0\)
\[
 (-1)^{m+1}R_m(t)=\int_0^t(-1)^{m+1}R_m'(s)\,ds
   =\int_0^t(-1)^mR_{m-1}(s)\,ds>0,
\]
because the integrand is continuous and positive on \((0,t]\). With
\(m=2n-1\) this says \(e^{-t}>\sum_{k=0}^{2n-1}(-t)^k/k!\), and with
\(m=2n\) it says \(e^{-t}<\sum_{k=0}^{2n}(-t)^k/k!\). This is (6.8).

For \(\lambda\geq0\) and \(\theta\geq0\), put \(t=\lambda\theta\). At
\(t=0\) all three members of (6.8) equal \(1\), so the non-strict form holds
for every \(t\geq0\). Integrate it against \(F\{d\theta\}\). The polynomial
members are integrable because \(\mu_1,\ldots,\mu_{2n}\) are finite, and
\(\int(\lambda\theta)^kF\{d\theta\}=\mu_k\lambda^k\). This is (6.7).

Suppose \(\sum_k\mu_k\lambda_0^k/k!<\infty\). For \(0\leq\lambda\leq\lambda_0\)
the terms \(\mu_k\lambda^k/k!\) are dominated by \(\mu_k\lambda_0^k/k!\), so
the alternating series in (6.9) converges absolutely, say to \(S\), and its
terms tend to zero. Let \(S_m\) be its partial sums. By (6.7),
\(S_{2n-1}\leq\varphi(\lambda)\leq S_{2n}\) for all \(n\geq1\), and both
bounds tend to \(S\). Hence \(\varphi(\lambda)=S\).
\end{proof}
\end{annexitem}

\begin{annexitem}{annex:p15-w02b-p15-012-laplace-moment-determinacy}{thm:p15-w02b-p15-012-laplace-moment-determinacy}{Real-variable propagation of the Laplace transform}
\begin{proof}
Let \(\varphi\) and \(\psi\) be the transforms of \(F\) and \(G\), and let
\(K=\sum_k\mu_k\lambda_0^k/k!<\infty\). For \(a\geq0\) put
\(\varphi^{[k]}(a)=(-1)^k\int_0^\infty e^{-a\theta}\theta^kF\{d\theta\}\),
and define \(\psi^{[k]}(a)\) from \(G\) in the same way. By
\cref{thm:p15-w02b-p15-012-laplace-derivatives},
\(\varphi^{[k]}(a)=\varphi^{(k)}(a)\) when \(a>0\).

\emph{Local expansion.} Let \(a\geq0\) and \(|\eta|\leq\lambda_0\) with
\(a+\eta\geq0\). Then
\(e^{-(a+\eta)\theta}=\sum_ke^{-a\theta}(-\eta\theta)^k/k!\) for every
\(\theta\geq0\), and
\[
 \sum_k\int_0^\infty e^{-a\theta}\frac{|\eta|^k\theta^k}{k!}\,F\{d\theta\}
    \leq\sum_k\frac{\mu_k\lambda_0^k}{k!}=K<\infty .
\]
By Tonelli's theorem applied to the absolute values, and then dominated
convergence for the partial sums, the series may be integrated term by term:
\[
 \varphi(a+\eta)=\sum_{k=0}^\infty\frac{\eta^k}{k!}\varphi^{[k]}(a),
\]
and likewise for \(\psi\), since \(G\) has the same moments.

\emph{Propagation.} At \(a=0\),
\(\varphi^{[k]}(0)=(-1)^k\mu_k=\psi^{[k]}(0)\), so the local expansion gives
\(\varphi=\psi\) on \([0,\lambda_0]\). Suppose \(\varphi=\psi\) on
\([0,b]\) for some \(b\geq\lambda_0\), and put \(a=b-\lambda_0/2\). Then
\(a\geq\lambda_0/2>0\) is an interior point of \([0,b]\). The derivatives of
every order at \(a\) are determined by the values on a neighborhood of \(a\)
inside \([0,b]\), where the two functions agree, so
\(\varphi^{[k]}(a)=\varphi^{(k)}(a)=\psi^{(k)}(a)=\psi^{[k]}(a)\) for all
\(k\). The local expansion with \(0\leq\eta\leq\lambda_0\) then gives
\(\varphi=\psi\) on \([a,a+\lambda_0]=[b-\lambda_0/2,b+\lambda_0/2]\), hence on
\([0,b+\lambda_0/2]\). By induction \(\varphi=\psi\) on
\([0,\lambda_0+j\lambda_0/2]\) for every \(j\), that is, on \([0,\infty)\).

\emph{Identification.} The two transforms agree for every \(\lambda>0\), so
\cref{thm:p15-w02b-p15-012-laplace-inversion} gives \(F=G\).
\end{proof}
\end{annexitem}

\begin{annexitem}{annex:p15-w02b-p15-012-laplace-examples}{thm:p15-w02b-p15-012-laplace-examples}{The exponential law, an atom and the lognormal moments}
\begin{proof}
(1) For \(\lambda\geq0\),
\(\varphi(\lambda)=\int_0^\infty e^{-\lambda\theta}e^{-\theta}\,d\theta=1/(1+\lambda)\).
Integration by parts gives
\(\int_0^\infty\theta^ke^{-\theta}\,d\theta=k\int_0^\infty\theta^{k-1}e^{-\theta}\,d\theta\),
so \(\mu_k=k!\) by induction. The series (6.9) is
\(\sum_k(-\lambda)^k=1/(1+\lambda)\) for \(0\leq\lambda<1\), and
\(\sum_k\mu_k\lambda_0^k/k!=\sum_k\lambda_0^k<\infty\) for
\(0<\lambda_0<1\), so \cref{thm:p15-w02b-p15-012-laplace-moment-determinacy}
applies. Differentiating gives
\(\varphi^{(k)}(\lambda)=(-1)^kk!(1+\lambda)^{-k-1}\), so
\((-1)^k\lambda^k\varphi^{(k)}(\lambda)/k!=q^k/(1+\lambda)\) with
\(q=\lambda/(1+\lambda)\). For \(x\geq0\) let \(K=\lfloor\lambda x\rfloor\).
Since \(1-q=1/(1+\lambda)\),
\[
 \sum_{k=0}^K\frac{q^k}{1+\lambda}=(1-q)\frac{1-q^{K+1}}{1-q}=1-q^{K+1}.
\]
Now \(q^{K+1}=\exp\{-(K+1)\log(1+1/\lambda)\}\). Here
\(\lambda x<K+1\leq\lambda x+1\), and
\(1/\lambda-1/(2\lambda^2)\leq\log(1+1/\lambda)\leq1/\lambda\) for
\(\lambda>0\), so \((K+1)\log(1+1/\lambda)\to x\) and the left side of (6.4)
tends to \(1-e^{-x}=F(x)\).

(2) For \(F=\delta_{x_0}\), \(\varphi(\lambda)=e^{-\lambda x_0}\) and
\((-1)^k\lambda^k\varphi^{(k)}(\lambda)/k!=e^{-\lambda x_0}(\lambda x_0)^k/k!\),
so the left side of (6.4) at \(x=x_0\) is
\(\mathbb P\{N_{\lambda x_0}\leq\lambda x_0\}\). At \(\lambda_m=m/x_0\) this
is \(\mathbb P\{N_m\leq m\}\). For independent \(N_a\) and \(N_b\),
\[
 \mathbb P\{N_a+N_b=k\}=\sum_{j=0}^ke^{-a}\frac{a^j}{j!}e^{-b}\frac{b^{k-j}}{(k-j)!}
   =e^{-(a+b)}\frac{(a+b)^k}{k!}
\]
by the binomial theorem, so \(N_m\) has the law of a sum \(S_m\) of \(m\)
independent Poisson variables of mean and variance \(1\). By
\cref{thm:p15-clt-iid}, \((S_m-m)/\sqrt m\) converges in distribution to a
standard normal variable \(\chi\). The point \(0\) is a continuity point of
the law of \(\chi\), so
\(\mathbb P\{N_m\leq m\}=\mathbb P\{(S_m-m)/\sqrt m\leq0\}\to
\mathbb P\{\chi\leq0\}=1/2\), whereas \(F(x_0)=1\).

(3) By \cref{thm:p15-feller-lognormal-moments}, the densities \(f_a\),
\(-1\leq a\leq1\), vanish on \((-\infty,0]\), are distinct, and all have the
moments \(\mu_k=e^{k^2/2}\). For \(\lambda>0\), \(k!\leq k^k\) gives
\[
 \frac{\mu_k\lambda^k}{k!}\geq\exp\Bigl\{k\Bigl(\frac k2+\log\lambda-\log k\Bigr)\Bigr\},
\]
and \(k/2-\log k\to\infty\), so the terms are unbounded and the series
diverges.
\end{proof}
\end{annexitem}

\begin{annexitem}{annex:p15-w02b-p15-012-exponential-remainder}{thm:p15-w02b-p15-012-exponential-remainder}{Feller's remainder estimate for the exponential}
\begin{proof}
The function \(x\mapsto e^{ix}\) has derivative \(ie^{ix}\), so
\(\rho_1(t)=e^{it}-1=i\int_0^te^{ix}\,dx\). For \(n>1\),
\[
 \rho_n'(t)=ie^{it}-\sum_{k=1}^{n-1}\frac{i(it)^{k-1}}{(k-1)!}
   =i\Bigl(e^{it}-\sum_{j=0}^{n-2}\frac{(it)^j}{j!}\Bigr)=i\rho_{n-1}(t),
\]
and \(\rho_n(0)=0\), so \(\rho_n(t)=i\int_0^t\rho_{n-1}(x)\,dx\).
For \(t>0\), \(|\rho_1(t)|\leq\int_0^t1\,dx=t\). If
\(|\rho_{n-1}(x)|\leq x^{n-1}/(n-1)!\) for \(x>0\), then
\(|\rho_n(t)|\leq\int_0^tx^{n-1}/(n-1)!\,dx=t^n/n!\). This is (4.1) for
\(t>0\). At \(t=0\) both sides vanish. The complex conjugate of
\(\rho_n(t)\) is \(e^{-it}-\sum_{k<n}(-it)^k/k!=\rho_n(-t)\), so
\(|\rho_n(-t)|=|\rho_n(t)|\leq|t|^n/n!\).
\end{proof}
\end{annexitem}

\begin{annexitem}{annex:p15-w02b-p15-012-xv4-derivatives}{thm:p15-w02b-p15-012-xv4-derivatives}{Feller's derivative formula by dominated convergence}
\begin{proof}
For \(0\leq j\leq n\), \(|x|^j\leq1+|x|^n\), so \(M_j\leq1+M_n<\infty\).
The case \(n=0\) is the definition. Suppose that for some \(k<n\),
\(\varphi^{(k)}(\zeta)=i^k\int e^{i\zeta x}x^kF\{dx\}\) for every \(\zeta\).
For real \(h\neq0\),
\[
 \frac{\varphi^{(k)}(\zeta+h)-\varphi^{(k)}(\zeta)}h
   =i^k\int e^{i\zeta x}x^k\frac{e^{ihx}-1}h\,F\{dx\}.
\]
By \cref{thm:p15-w02b-p15-012-exponential-remainder} with \(n=1\),
\(|e^{ihx}-1|\leq|hx|\), so the integrand is dominated by \(|x|^{k+1}\),
which is integrable since \(M_{k+1}<\infty\). As \(h\to0\) along any sequence,
\((e^{ihx}-1)/h\to ix\). Dominated convergence gives
\(\varphi^{(k+1)}(\zeta)=i^{k+1}\int e^{i\zeta x}x^{k+1}F\{dx\}\). This proves
(4.5) by induction. If \(\zeta_j\to\zeta\), the integrands
\(e^{i\zeta_jx}x^n\) converge pointwise and are dominated by \(|x|^n\), so
\(\varphi^{(n)}\) is continuous. At \(\zeta=0\) the formula gives
\(i^nm_n\).
\end{proof}
\end{annexitem}

\begin{annexitem}{annex:p15-w02b-p15-012-taylor-bound}{thm:p15-w02b-p15-012-taylor-bound}{The Taylor bound for characteristic functions}
\begin{proof}
Since \(|e^{i\zeta x}|=1\), the left side of (4.13) equals
\(|\rho_n(tx)|\), which is at most \(|tx|^n/n!\) by (4.1).

Let \(M_n<\infty\). By \cref{thm:p15-w02b-p15-012-xv4-derivatives},
\(\varphi^{(k)}(\zeta)=\int e^{i\zeta x}(ix)^kF\{dx\}\) for \(k\leq n-1\), so
\[
 \varphi(\zeta+t)-\sum_{k=0}^{n-1}\frac{t^k}{k!}\varphi^{(k)}(\zeta)
   =\int e^{i\zeta x}\Bigl(e^{itx}-\sum_{k=0}^{n-1}\frac{(itx)^k}{k!}\Bigr)F\{dx\}.
\]
By (4.13) the modulus of the integrand is at most \(|tx|^n/n!\), and
integration gives the bound \(M_n|t|^n/n!\). For \(F=\delta_0\) and
\(n\geq1\) both sides of (4.14) are zero, so the inequality cannot be strict
in general.
\end{proof}
\end{annexitem}

\begin{annexitem}{annex:p15-w02b-p15-012-xv4-determinacy}{thm:p15-w02b-p15-012-xv4-determinacy}{Feller's analytic determinacy proof, with the propagation written out}
\begin{proof}
\emph{A factorial bound.} For \(n\geq1\),
\(e^n=\sum_jn^j/j!\geq n^n/n!\), so \(n!\geq(n/e)^n\).

\emph{A uniform moment bound.} Let \(\Lambda>\lambda\) with \(\Lambda>0\).
By (4.15) there is \(n_0\) with \(M_n^{1/n}\leq\Lambda n\) for \(n\geq n_0\).
Put \(C=\max\{1,M_n/(\Lambda n)^n:1\leq n<n_0\}\). Then
\(M_n\leq C(\Lambda n)^n\) for every \(n\geq1\).

\emph{The Taylor series of \(\varphi\).} Let \(\zeta,t\) be real with
\(e\lambda|t|<1\); for \(t=0\) there is nothing to prove. Choose
\(\Lambda\) with \(\max(\lambda,0)<\Lambda<1/(e|t|)\). By (4.14) and the two
bounds,
\[
 \Bigl|\varphi(\zeta+t)-\sum_{k=0}^{n-1}\frac{t^k}{k!}\varphi^{(k)}(\zeta)\Bigr|
   \leq C(\Lambda n)^n\frac{e^n|t|^n}{n^n}=C(e\Lambda|t|)^n\longrightarrow0 .
\]
This is the series assertion. When \(\lambda=0\) every \(t\) is admissible.
If \(|t|<1/(3\lambda)\) then \(e\lambda|t|<3\lambda|t|<1\), which gives
Feller's radius.

\emph{A competing law satisfies the same bound.} Let \(G\) have the same
moments \(m_n\), with absolute moments \(M_n^G\). Even absolute moments are
moments, so \(M_{2k}^G=m_{2k}=M_{2k}\). By the Cauchy--Schwarz inequality
applied to \(|x|^k\cdot|x|^{k+1}\),
\[
 M_{2k+1}^G\leq(M_{2k}M_{2k+2})^{1/2}
   \leq C\bigl(\Lambda(2k+2)\bigr)^{(2k+2k+2)/2}=C\bigl(\Lambda(2k+2)\bigr)^{2k+1},
\]
using \((\Lambda\cdot2k)^{2k}\leq(\Lambda(2k+2))^{2k}\), with the convention
\(0^0=1\) at \(k=0\). Hence \(M_n^G\leq C(\Lambda(n+1))^n\) for every
\(n\geq1\), and
\[
 M_n^G\frac{|t|^n}{n!}\leq C(e\Lambda|t|)^n\Bigl(1+\frac1n\Bigr)^n
   \leq Ce\,(e\Lambda|t|)^n .
\]

\emph{Propagation along the real axis.} Fix \(\Lambda>\max(\lambda,0)\) and
\(t_0>0\) with \(e\Lambda t_0<1\). Let \(\varphi\) and \(\psi\) be the
characteristic functions of \(F\) and \(G\), and \(D=\varphi-\psi\). By
(4.14) for each law and the two bounds, for every real \(\zeta\) and
\(|t|\leq t_0\),
\[
 D(\zeta+t)=\sum_{k=0}^\infty\frac{t^k}{k!}D^{(k)}(\zeta).
\]
Let \(S\) be the set of real \(\zeta\) at which \(D^{(k)}(\zeta)=0\) for every
\(k\geq0\). Each \(D^{(k)}\) is continuous by
\cref{thm:p15-w02b-p15-012-xv4-derivatives}, so \(S\) is closed. By (4.5),
\(\varphi^{(k)}(0)=i^km_k=\psi^{(k)}(0)\), so \(0\in S\). If \(\zeta\in S\),
the expansion gives \(D=0\) on \([\zeta-t_0,\zeta+t_0]\); all derivatives of
\(D\) vanish on the open interval \((\zeta-t_0,\zeta+t_0)\), and by closedness
\([\zeta-t_0,\zeta+t_0]\subseteq S\). Starting from \(0\) and applying this to
the endpoints \(\pm jt_0\), induction gives \([-jt_0,jt_0]\subseteq S\) for all
\(j\). Hence \(D\equiv0\), so \(\varphi=\psi\) on the line, and
\cref{thm:p15-characteristic-inversion-uniqueness} gives \(F=G\).
\end{proof}
\end{annexitem}

\begin{annexitem}{annex:p15-w02b-p15-012-xv4-bridge}{thm:p15-w02b-p15-012-xv4-bridge}{From the moment power series to (4.15) and to Carleman's condition}
\begin{proof}
Let \(C'=\sum_nm_{2n}s^{2n}/(2n)!\). The term \(n=0\) is \(m_0=1\), so
\(1\leq C'<\infty\), and every term is at most \(C'\):
\[
 M_{2n}=m_{2n}\leq C'(2n)!\,s^{-2n}\leq C'(2n)^{2n}s^{-2n}.
\]
Hence \(M_{2n}^{1/(2n)}/(2n)\leq C'^{1/(2n)}/s\to1/s\). For odd orders,
Cauchy--Schwarz gives
\[
 M_{2n+1}\leq(M_{2n}M_{2n+2})^{1/2}\leq C'(2n+2)^{2n+1}s^{-(2n+1)},
\]
so \(M_{2n+1}^{1/(2n+1)}/(2n+1)\leq C'^{1/(2n+1)}(2n+2)/((2n+1)s)\to1/s\).
Thus the upper limit in (4.15) is at most \(1/s\), and
\cref{thm:p15-w02b-p15-012-xv4-determinacy} gives determinacy.
Putting \(t=s^2\) turns the hypothesis into convergence of
\(\sum_nm_{2n}t^n/(2n)!\), and conversely.
For a law on \([0,\infty)\), the terms \(\mu_{2n}\lambda_0^{2n}/(2n)!\) form a
subseries of the convergent nonnegative series \(\sum_k\mu_k\lambda_0^k/k!\),
so the hypothesis holds with \(s=\lambda_0\).

Assume (4.15) and let \(\Lambda>\max(\lambda,0)\). For \(n\geq n_0\),
\(M_n^{1/n}\leq\Lambda n\), so \(M_n^{-1/n}\geq1/(\Lambda n)\), the value
\(+\infty\) being allowed when \(M_n=0\). The harmonic series diverges, so
\(\sum_nM_n^{-1/n}=\infty\). The even-order subseries
\(\sum_km_{2k}^{-1/(2k)}\geq\sum_{k\geq n_0}1/(2\Lambda k)\) also diverges,
which is the form used in \cref{thm:p15-criteria-carleman}.
\end{proof}
\end{annexitem}

\begin{annexitem}{annex:p15-w02b-p15-012-xv4-examples}{thm:p15-w02b-p15-012-xv4-examples}{Normal, exponential and lognormal growth}
\begin{proof}
(1) By \cref{thm:p15-criteria-moment-identities},
\(M_{2k}=(2k-1)!!\), a product of \(k\) odd numbers each at most \(2k\), so
\(M_{2k}\leq(2k)^k\) and \(M_{2k}^{1/(2k)}/(2k)\leq(2k)^{-1/2}\to0\). By
Cauchy--Schwarz,
\(M_{2k+1}\leq((2k)^k(2k+2)^{k+1})^{1/2}\leq(2k+2)^{(2k+1)/2}\), so
\(M_{2k+1}^{1/(2k+1)}/(2k+1)\leq(2k+2)^{1/2}/(2k+1)\to0\). Thus
\(\lambda=0\), and \cref{thm:p15-w02b-p15-012-xv4-determinacy} gives the
expansion for every real increment.

(2) The exponential law is concentrated on \([0,\infty)\), so
\(M_n=\mu_n=n!\) by the computation in
\cref{thm:p15-w02b-p15-012-laplace-examples}. Since \(n!\leq n^n\),
\(M_n^{1/n}/n\leq1\), and the upper limit is at most \(1\).

(3) The lognormal laws are concentrated on \((0,\infty)\), so
\(M_n=e^{n^2/2}\) and \(M_n^{1/n}/n=e^{n/2}/n\to\infty\). Condition (4.15)
fails; by \cref{thm:p15-feller-lognormal-moments} these moments are shared by
distinct laws, so no sufficient criterion can hold for them.
\end{proof}
\end{annexitem}
